% ============================================================
%  VARAKUMAR G  —  Résumé (Sector base: Detection Engineering / SOC)
%  Derived from ../varakumar_resume.tex — SAME facts/numbers/entities, only the
%  lens (which aspect of each fact is foregrounded) changes. See ../resumes/facts.json
%  for the fact ledger this file is built from. Do not add any claim not already
%  present in the master. PLAN.md §9 2026-08-24.
%  Engine : tectonic (pdfTeX-compatible)
% ============================================================
\documentclass[9pt,a4paper]{extarticle}

% ── Fonts & Encoding ────────────────────────────────────────
\usepackage[sfdefault]{inter} % Inter — modern geometric sans-serif
\usepackage{unicode-math}     % OpenType math — scales cleanly to any pt size

% ── Page Layout ─────────────────────────────────────────────
\usepackage[top=0.45in, bottom=0.45in, left=0.50in, right=0.50in]{geometry}
\setlength{\parindent}{0pt}
\setlength{\parskip}{0pt}
\pagestyle{empty}

% ── Tight Bullet Lists ──────────────────────────────────────
\usepackage{enumitem}
\setlist[itemize]{
  leftmargin = 1.2em,
  label      = \textbullet,
  itemsep    = 0pt,
  parsep     = 0pt,
  topsep     = 1pt,
  partopsep  = 0pt
}

% ── Links  (real text in PDF, no colour boxes — ATS safe) ───
\usepackage[hidelinks]{hyperref}
\urlstyle{same}

% ── Micro-typography ────────────────────────────────────────
\usepackage{microtype}

% ── Font size: 8.5pt body (slightly smaller than 9pt default) ──
\usepackage{fontsize}
\changefontsize[10.7pt]{8.5pt}

% ── Macros ──────────────────────────────────────────────────
\newcommand{\resumeSection}[1]{%
  \vspace{5pt}%
  \noindent{\large\bfseries\MakeUppercase{#1}}\par%
  \vspace{-7pt}%
  \noindent\rule{\textwidth}{0.5pt}%
  \par\vspace{0pt}%
}
\newcommand{\techrow}[2]{\noindent\textbf{#1} #2\par}


% ============================================================
\begin{document}
% ============================================================

% ── HEADER ──────────────────────────────────────────────────
\vspace*{-8pt}
\begin{center}
  {\LARGE\bfseries Varakumar G}\\[2pt]
  \href{mailto:varakumar7000@gmail.com}{varakumar7000@gmail.com}
  ~|~ +91\,6309856263
  ~|~ \href{https://www.linkedin.com/in/varakumar01}{linkedin.com/in/varakumar01}\\[1pt]
  \href{https://github.com/varakumar01}{GitHub}
  ~|~ \href{https://gitlab.com/Tony_01}{GitLab}
  ~|~ \href{https://portfolio.dragnux.com}{Portfolio}
  ~|~ \href{https://drive.google.com/file/d/17c5RtAgWn5TdRQUxnDIvLQogPoUoxgO_/view?pli=1}{Letter of Recommendation}
\end{center}
\vspace{-1pt}\noindent\rule{\textwidth}{0.7pt}

% ── PROFESSIONAL SUMMARY ────────────────────────────────────
\resumeSection{Professional Summary}

Security Engineer with 2+ years in \textbf{detection engineering} and vulnerability research — authoring custom \textbf{NASL} detection plugins for \textbf{OT/SCADA CVEs} and growing coverage from \textbf{~400 to 3,200+} within a year, a \textbf{700\%} increase, with automated feed updates as new CVEs land. Builds \textbf{MITRE ATT\&CK}-aligned coverage across OT and network protocols (Modbus, BACnet, DNP3, PROFINET, IPsec, Kerberos) plus cloud posture signals via \textbf{CloudSploit}, and validates every check offensively — \textbf{Active Directory} attack paths and web app exploitation — so detections fire on real attack behavior, not stale signatures.

Automated plugin development for OT security to keep pace with the volume, and I contribute to \textbf{open-source} security testing projects because I want the tools I depend on to stay secure.

% ── TECHNICAL ENVIRONMENT ───────────────────────────────────
\resumeSection{Technical Skills}

\techrow{Security Domains:}{Detection Engineering, Vulnerability Research, Vulnerability Management, ICS/OT Security, Cloud Security}
\techrow{Detection \& Scanning:}{OpenVAS, NASL, Nessus, Qualys, CloudSploit, NVD, CVSS, Custom Scripts}
\techrow{Frameworks \& Standards:}{MITRE ATT\&CK, OWASP Top 10, NIST CSF, CIS Benchmarks, CVSS, PCI-DSS, ISO 27001}
\techrow{Network \& Security Protocols:}{IPsec, Kerberos, LDAP, TLS/SSL, DNS, DHCP, SNMP, SMB, RDP, SSH, HTTP/HTTPS}
\techrow{OT/ICS \& Industrial Protocols:}{Modbus, BACnet, EtherNet/IP, DNP3, PROFINET}
\techrow{Cloud Platforms:}{AWS, Microsoft Azure, Google Cloud Platform (GCP), Oracle Cloud}
\techrow{Exploitation \& Security Testing:}{Metasploit, Burp Suite, Nmap, Wireshark, Postman, crAPI, VAmPI}
\techrow{Programming \& Platforms:}{C, Python, Node.js, Bash, MongoDB, MySQL, Kali Linux, Arch Linux, WireGuard}

% ── PROFESSIONAL EXPERIENCE ─────────────────────────────────
\resumeSection{Professional Experience}

\noindent\textbf{Holm Security AB} | Security Engineer \hfill \textbf{Randstad.in}
\begin{itemize}
  \item Authored custom \textbf{NASL} detection plugins for \textbf{OT/SCADA} vulnerabilities, growing
        \textbf{CVE coverage from ~400 to 3,200+} within a year — a \textbf{700\% increase} in the company's
        SCADA/OT detection coverage for clients.
  \item Built and tuned detection logic against OT products from \textbf{Siemens}, \textbf{Schneider Electric},
        \textbf{Rockwell Automation}, \textbf{ABB}, \textbf{Emerson}, and \textbf{Honeywell}, eliminating false
        positives and scoring findings with \textbf{CVSS} so alert volume tracked real exploitability.
  \item Managed the \textbf{CloudSploit platform end-to-end} across \textbf{AWS}, \textbf{Azure}, \textbf{GCP}, and
        \textbf{Oracle}, maintaining \textbf{CIS Benchmark} posture checks and contributing fixes upstream.
  \item Automated detection-content development using \textbf{Claude AI}, keeping the feed current as new CVEs
        were published, and re-validated coverage with clients after each release.
\end{itemize}

\vspace{3pt}
\noindent\textbf{Independent Security Assessment}
\begin{itemize}
  \item \textbf{Bill Manager (E-commerce/Billing App) — Independent Web Application Assessment:}
  \begin{itemize}
    \item Assessed and exploited a live billing application end-to-end: \textbf{broken access control}
          between user and admin roles, \textbf{insecure direct object reference (IDOR)} on
          invoices/billing records, \textbf{privilege escalation} via role/parameter tampering,
          \textbf{SQL injection} on login and search fields, \textbf{stored/reflected XSS}, an insecure
          password-reset flow, and weak \textbf{session/JWT} handling.
    \item Reported every finding with remediation guidance through a \textbf{GitHub Security Advisory}
          on the project repository.
  \end{itemize}

  \item \textbf{Active Directory Security Assessment:}
  \begin{itemize}
    \item Executed core AD attack paths end-to-end: \textbf{Kerberoasting}, \textbf{LLMNR/NBT-NS
          poisoning}, \textbf{pass-the-hash}, and \textbf{privilege escalation} via misconfigured \textbf{ACLs}.
    \item Building toward advanced AD attack-path expertise: \textbf{ADCS abuse (ESC1 to ESC8)},
          \textbf{delegation attacks} (unconstrained, constrained, resource-based),
          \textbf{Shadow Credentials}, and \textbf{BloodHound}-driven attack path analysis.
  \end{itemize}

  \item \textbf{Web \& API Security Testing (PortSwigger Web Security Academy):}
  \begin{itemize}
    \item Exploited \textbf{SQL injection}, \textbf{XSS}, \textbf{CSRF}, \textbf{authentication bypass},
          \textbf{access control}, \textbf{SSRF}, \textbf{XXE}, and \textbf{insecure deserialization} labs
          hands-on, carrying the same techniques into \textbf{API security testing} on \textbf{crAPI}/\textbf{VAmPI}.
  \end{itemize}
\end{itemize}

% ── PROJECTS ────────────────────────────────────────────────
\resumeSection{Projects}
\begin{itemize}
  \item \textbf{Subdomain \& Attack Surface Enumeration Tool}
        | Python, Requests, dnspython \\
        Turns a target scope into a prioritized attack surface: enumerates subdomains from public sources
        like \textbf{crt.sh}, \textbf{brute-forces} the remainder over \textbf{DNS}, then scans open ports and
        fingerprints each host's stack with \textbf{OpenVAS} plus custom backend-fingerprinting plugins.

  \item \textbf{Self-Maintained OpenVAS Scanner}
        | Nmap, NASL \\
        Maintain a fork of \textbf{OpenVAS} carrying custom \textbf{NASL} plugins for \textbf{CVE} checks the default
        feed doesn't cover, plus fingerprinting logic for newer software. AI-automated scripts keep the feed
        current as CVEs are published, so findings map to exploitable conditions, not stale signatures.

  \item \textbf{Local Lab Setup}
        \begin{itemize}
        \item \textbf{Active Directory lab:} Built a \textbf{domain controller} and client machines to run full
              attack chains hands-on — \textbf{Kerberoasting}, \textbf{LLMNR poisoning}, \textbf{privilege escalation}.
        \item \textbf{pfSense + Vulnerable VMs:} Home lab behind a \textbf{pfSense} firewall with intentionally
              vulnerable targets (\textbf{Metasploitable}, \textbf{DVWA}, \textbf{VulnHub}) covering the full
              recon-to-exploit flow.
        \end{itemize}
\end{itemize}

% ── CERTIFICATIONS ──────────────────────────────────────────
\resumeSection{Certifications}

\begin{itemize}
  \item Offensive Security Certified Professional (OSCP) | OffSec | \textit{In Progress}
  \item TCM Security: Practical Ethical Hacking \textbullet\ Windows \& Linux Privilege Escalation \textbullet\ OSINT Fundamentals \textbullet\ External Pentest Playbook
\end{itemize}

% ── EDUCATION ───────────────────────────────────────────────
\resumeSection{Education}

\noindent\textbf{Malla Reddy College of Engineering \& Technology}, Hyderabad, India%
\hfill 2018 -- 2022\\
Bachelor of Technology (B.Tech), Computer Science

% ============================================================
\end{document}
